\section*{Hoofdstuk \ref{ch:b3:statistical-thermo-applied} --- Toegepaste statistische thermodynamica}

\begin{solution}{exo:b3:statistical-thermo-applied:1}
$\frac32R + R + (3 \times 3 - 5)R = 6.5R = \qty{54.0}{J.K^{-1}.mol^{-1}}$. Bij
kamertemperatuur zijn translatie en rotatie klassiek, zijn de twee
strekvibraties bevroren en is de tweevoudig ontaarde buigvibratie, de
laagste vibratie, gedeeltelijk aangeslagen.
\end{solution}

\begin{solution}{exo:b3:statistical-thermo-applied:2}
$x = 1$: $\eu/(\eu - 1)^2 = 0.921$. $x = 10$: $100\eu^{10}/(\eu^{10} - 1)^2 \approx 100\eu^{-10} = \num{4.5e-3}$.
\end{solution}

\begin{solution}{exo:b3:statistical-thermo-applied:3}
\qty{20}{K}: $x_{\mathrm{para}} = 1/(1 + 3 \times 3\eu^{-2 \times 85.35/20}) = 1/(1 + 9 \times \num{1.96e-4}) = 0.998$.
\qty{77}{K}: even som $1 + 5\eu^{-6.651} = 1.0065$; oneven som $3(3\eu^{-2.217} + 7\eu^{-13.30}) = 0.9806$;
$x_{\mathrm{para}} = 1.0065/1.9871 = 0.51$.
\end{solution}

\begin{solution}{exo:b3:statistical-thermo-applied:4}
Voor $I = 1$ zijn er $3 \times 2 = 6$ symmetrische en $3 \times 1 = 3$
antisymmetrische spintoestanden; deuteronen zijn bosonen, dus de totale
golffunctie is symmetrisch: symmetrische spintoestanden bij even $J$
(ortho, gewicht 6), antisymmetrische bij oneven $J$ (para, gewicht
3). Orthodeuterium, dat $J = 0$ bevat, is stabiel bij 0 K; bij
kamertemperatuur is ortho:para = 2:1.
\end{solution}

\begin{solution}{exo:b3:statistical-thermo-applied:5}
$\Delta_rE_0 = 2 \times 1883.7 - 2170.3 - 1542.3 = \qty{54.8}{cm^{-1}}$; $\eu^{-54.8/207.2} = 0.768$, en
$4 \times 0.768 = 3.07$. De volledige waarde, 3.26, is 6\,\% hoger: de
translatie-, rotatie- en vibratiefactoren vallen alleen in de klassieke
limiet precies weg, en de rotatie van deze lichte moleculen is bij
\qty{298}{K} nog gedeeltelijk gekwantiseerd.
\end{solution}

\begin{solution}{exo:b3:statistical-thermo-applied:6}
\ce{N2O}: twee oriëntaties, $R\ln2 = \qty{5.76}{J.K^{-1}.mol^{-1}}$. \ce{CH3D}: vier, $R\ln4 =
\qty{11.53}{J.K^{-1}.mol^{-1}}$ (bovengrenzen, bereikt als de schikkingen
willekeurig zijn).
\end{solution}

\begin{solution}{exo:b3:statistical-thermo-applied:7}
$x = 797.6/300 = 2.659$: einsteinfunctie $x^2\eu^{-x}/(1 - \eu^{-x})^2 = 0.572$; $C_V/R = 3.072$,
$C_p = 4.072R = \qty{33.86}{J.K^{-1}.mol^{-1}}$, 0.4\,\% onder de tabel
(anharmoniciteit).
\end{solution}

\begin{solution}{exo:b3:statistical-thermo-applied:8}
$x = 1.029$: $x^2\eu^{-x}/(1 - \eu^{-x})^2 = 0.916$, dus \qty{7.62}{J.K^{-1}.mol^{-1}},
92\,\% van $R$: de vibratie van jood is bij kamertemperatuur bijna
klassiek.
\end{solution}

\begin{solution}{exo:b3:statistical-thermo-applied:9}
B is rijker aan \ce{^{18}O}. $\alpha_{\mathrm{A/B}} = (1 - 0.0100)/(1 + 0.0020) = 0.98802$.
\end{solution}

\begin{solution}{exo:b3:statistical-thermo-applied:10}
$\Delta\mathrm{ZPE} = \frac12(1 - 1/\sqrt2)(2000 - 3000) = \qty{-146}{cm^{-1}}$; $K = \eu^{146.4/207.2} = 2.0$.
Het deuterium gaat bij voorkeur naar X, de stijvere binding, waar zijn
besparing aan \omterm{def:b3:quantum-model-systems:oscillator}{nulpuntsenergie} het grootst is.
\end{solution}

\begin{solution}{exo:b3:statistical-thermo-applied:11}
$\Delta_rH^\circ = R\ln10\,(-1.766 + 3.392)/(1/900 - 1/1100) = 19.145 \times 1.626/\num{2.020e-4} =
\qty{154}{kJ/mol}$. De twee atomen dragen $2 \times \frac52RT$ aan
enthalpie; het molecuul $\frac52RT + RT$
(rotatie) plus zijn vibratie-energie $R\theta_{\mathrm V}/(\eu^{\theta_{\mathrm V}/T} - 1)$, iets minder dan $RT$:
het verschil, ongeveer \qty{5}{kJ/mol} bij \qty{1000}{K}, komt bij $D_0$.
\end{solution}

\begin{solution}{exo:b3:statistical-thermo-applied:12}
Niveaus 0 ($g = 1$) en $6hcB$ ($g = 5$): $q = 1 + 5\eu^{-x}$, $\langle\varepsilon\rangle/kT = 5x\eu^{-x}/q$, $\langle\varepsilon^2\rangle/(kT)^2 =
5x^2\eu^{-x}/q$, en $C/R = \langle\varepsilon^2\rangle/(kT)^2 - (\langle\varepsilon\rangle/kT)^2 = 5x^2\eu^{-x}/(1 + 5\eu^{-x})^2$. Bij
\qty{100}{K} is $x = 5.121$: $C/R = 0.783/1.061 = 0.738$.
\end{solution}

\begin{solution}{pb:b3:statistical-thermo-applied:1}
\textbf{1.} $2 \times 2 = 4$: drie symmetrische (het triplet), één
antisymmetrische (het singulet).
\textbf{2.} Protonen zijn fermionen, dus de totale golffunctie verandert
van teken bij verwisseling; de rotatiefunctie geeft $(-1)^J$: even $J$
hoort bij het antisymmetrische singulet (para), oneven $J$ bij het
symmetrische triplet (ortho).
\textbf{3.} Even niveaus 1, oneven niveaus 3.
\textbf{4.} Bij \qty{300}{K} ($T = 3.5\,\theta_{\mathrm R}$) zijn de even en de
oneven rotatiesom bijna gelijk; gewogen met 1 en 3 geven ze 1:3
(parafractie 0.2507).
\textbf{5.} $F(1) = 2 \times 59.322 - 4 \times 0.0471 = \qty{118.46}{cm^{-1}} = \qty{1.417}{kJ/mol}$.
\textbf{6.} $1/(1 + 9\eu^{-170.4/20.37}) = 1/(1 + 9 \times \num{2.3e-4}) = 0.998$.
\textbf{7.} Met alleen $J = 0$ en 1 betekent half para $9\eu^{-2\theta_{\mathrm R}/T} = 1$, $T = 2\theta_{\mathrm R}/\ln9 =
0.91\,\theta_{\mathrm R} = \qty{78}{K}$: het is een wedstrijd tussen het
gewicht 9 van $J = 1$ en zijn boltzmannfactor.
\textbf{8.} 0.75. Zonder katalysator vraagt de omzetting dagen, het
vloeibaar maken uren.
\textbf{9.} $(0.75 - 0.002) \times 1.417 = \qty{1.060}{kJ/mol}$.
\textbf{10.} $1.060/\num{2.016e-3} = \qty{526}{kJ/kg}$.
\textbf{11.} $0.905/\num{2.016e-3} = \qty{449}{kJ/kg}$.
\textbf{12.} De omzettingswarmte is 1.17 keer de verdampingsenthalpie.
\textbf{13.} Ze ontstaat in de vloeistof zelf en komt niet via de wanden
binnen.
\textbf{14.} $x_{\mathrm o}(t) = x_{\mathrm{eq}} + (0.75 - x_{\mathrm{eq}})\eu^{-kt}$, met $x_{\mathrm{eq}} = 0.002$.
\textbf{15.} $0.002 + 0.748\eu^{-0.274} = 0.571$.
\textbf{16.} $(0.75 - 0.571) \times 1.417 = \qty{0.254}{kJ/mol}$.
\textbf{17.} $0.254/0.905 = 0.28$: 28\,\% van de tank op één dag.
\textbf{18.} Na \qty{168}{h} is $x_{\mathrm o} = 0.112$, warmte
\qty{0.904}{kJ/mol}, gelijk aan de verdampingsenthalpie: in dit model is
de tank leeg. De schatting overschat het verlies omdat de verdampte
moleculen hun niet-omgezette \omterm{def:b3:statistical-thermo-applied:spin-isomers}{orthowaterstof} meenemen; het werkelijke
verlies is kleiner, maar nog steeds groot.
\textbf{19.} $t_{1/2} = \ln2/0.0114 = \qty{61}{h}$.
\textbf{20.} In de liquefactor, op katalysatorbedden bij opeenvolgende
temperaturen, zodat de koelmachine de omzettingswarmte afvoert voordat de
vloeistof opgeslagen wordt.
\textbf{21.} Ja: de parafractie in evenwicht bij \qty{300}{K} is 0.2507.
\textbf{22.} In normale waterstof kunnen de twee vormen niet in elkaar
overgaan, zodat de rotatie van elk ervan bevriest op haar eigen ladder van
niveaus; in waterstof in evenwicht zet een deel van de toegevoerde warmte
para om in ortho.
\textbf{23.} Ongeveer 3.57 rond \qty{50}{K}: naast het aanslaan van de
rotatie zet de warmte moleculen om van $J = 0$ naar $J = 1$,
\qty{118}{cm^{-1}} hoger, een reactie met een positieve enthalpie.
\textbf{24.} 1.50: de rotatie is in beide vormen bevroren (para in $J = 0$,
ortho in $J = 1$).
\textbf{25.} \textbf{Omzettingswarmte / verdampingsenthalpie $\approx 1.2$
voor normale waterstof}: de omzetting alleen al zou de hele tank kunnen
doen verdampen, en daarom wordt waterstof vóór de opslag naar para
omgezet.
\end{solution}
